\chapter{Jeux, enchères et information}\label{ch:qm:games-auctions-and-information}

En novembre 1998, le Trésor américain a fait passer toutes les adjudications de titres négociables à un prix de compensation unique~: tous les adjudicataires paient le rendement le plus élevé accepté, au lieu de payer chacun son propre prix offert.
Il avait mené depuis septembre 1992 une expérience portant sur des bons du Trésor à deux et cinq ans, et son étude propre de cette expérience a montré que les dealers formulaient des offres plus agressives sous le prix unique, que les attributions étaient moins
concentrées, et que l'écart de rendement de l'adjudication par rapport au marché when-issued a à peu près été divisé par deux pour les bons à deux ans (de 0.41 à 0.22 point de base), une différence qui n'était pas statistiquement
significative, car les résultats fluctuaient davantage d'une adjudication à l'autre. Un dealer formule des offres différemment selon les deux formats. La théorie des jeux et des enchères dit comment, et si le Trésor aurait dû
s'attendre à lever davantage. Ce chapitre développe cette théorie et se termine par la théorie de l'information d'un parieur qui achète un tuyau.

\section{Jeux et équilibre}

\begin{definition}[Jeu sous forme normale, équilibre de Nash, stratégie mixte, jeu à somme nulle]\label{def:qm:games-auctions-and-information:game}
Un \emph{jeu sous forme normale}\index{jeu sous forme normale} est un ensemble de joueurs, un ensemble de stratégies pour chacun et un gain pour chaque joueur pour chaque profil de stratégies. Une \emph{stratégie mixte}\index{stratégie mixte} est une
distribution de probabilité sur les stratégies d'un joueur. Un \emph{équilibre de Nash}\index{equilibre@équilibre de Nash} est un profil de stratégies (mixtes) dans lequel aucun joueur ne gagne à dévier seul. Un \emph{jeu à somme
nulle}\index{jeu à somme nulle} est un jeu à deux joueurs dont la somme des gains vaut zéro~; avec une matrice de gains $A$ pour le joueur en ligne, les stratégies mixtes $x$ et $y$ lui donnent $x^\top Ay$.
\end{definition}

\begin{theorem}[Existence d'équilibres~; minimax]\label{thm:qm:games-auctions-and-information:minimax}
Tout jeu fini admet un \omterm{def:qm:games-auctions-and-information:game}{équilibre de Nash} en \omterm{def:qm:games-auctions-and-information:game}{stratégies mixtes} (Nash, 1950~; admis). Dans un jeu fini à somme nulle, $\max_x\min_yx^\top Ay = \min_y\max_xx^\top Ay$ (von Neumann, 1928)~; la valeur commune est la valeur
du jeu, et les stratégies optimales forment les équilibres.
\end{theorem}

\begin{proof}[Idée]
L'existence découle du théorème du point fixe de Kakutani appliqué à la correspondance de meilleure réponse. Pour le théorème du minimax, le problème du joueur en ligne $\max\{v : A^\top x \ge v\mathbf 1, \mathbf 1^\top x = 1, x \ge
0\}$ est un \omterm{def:qm:convex-optimisation:classes}{programme linéaire} dont le dual est le problème du joueur en colonne, et la \omterm{def:qm:convex-optimisation:strong}{dualité forte} (chapitre~23) égalise leurs valeurs.
\end{proof}

Dans le pierre--papier--ciseaux pondéré (la pierre bat les ciseaux pour 1, les ciseaux battent le papier pour 1, le papier bat la pierre pour 2), l'équilibre n'est pas uniforme~: chaque joueur joue pierre, papier et ciseaux avec les probabilités
$(\frac14, \frac12, \frac14)$, et la valeur est nulle. Deux joueurs qui appliquent chacun la méthode des poids multiplicatifs pendant 20\,000 tours, en renforçant les actions qui auraient rapporté, donnent en moyenne $(0.250, 0.502, 0.248)$ avec un
écart de dualité de 0.005~: la dynamique d'apprentissage trouve l'équilibre qu'aucun des deux joueurs n'a calculé. Les \omterm{def:qm:games-auctions-and-information:game}{équilibres de Nash} des jeux à somme non nulle sont plus difficiles à calculer et ne sont pas nécessairement uniques~; un desk de trading les rencontre dans
les jeux de cotation (le livre~2) et les enchères ci-dessous.

\section{Jeux bayésiens}

\begin{definition}[Jeu bayésien, équilibre de Bayes--Nash]\label{def:qm:games-auctions-and-information:bayes}
Un \emph{jeu bayésien}\index{jeu bayésien} (Harsanyi, 1967--68) ajoute à un jeu un type pour chaque joueur, tiré d'un a priori commun et connu de son seul détenteur~; une stratégie associe des types à des actions. Un
\emph{équilibre de Bayes--Nash}\index{equilibre@équilibre de Bayes--Nash} est un profil de stratégies dans lequel l'action de chaque type maximise son gain espéré compte tenu de ses croyances sur les types des autres et sur leurs stratégies.
\end{definition}

\begin{definition}[Valeurs privées et valeurs communes, enchères au premier et au second prix]\label{def:qm:games-auctions-and-information:auctions}
Dans une \emph{enchère à valeurs privées}\index{enchère à valeurs privées} chaque enchérisseur connaît sa propre valeur pour l'objet~; dans une \emph{enchère à valeur commune}\index{enchère à valeur commune} la valeur est la même pour tous et chacun
n'en observe qu'un signal. Dans une \emph{enchère au premier prix}\index{enchère au premier prix} les offres scellées sont ouvertes et l'enchérisseur le plus haut paie son offre~; dans une \emph{enchère au second prix}\index{enchère au second prix}
(Vickrey, 1961), l'enchérisseur le plus haut paie la deuxième offre la plus élevée. La \emph{sous-enchère stratégique}\index{sous-enchère stratégique} (bid shading) est le fait d'enchérir en dessous de sa propre valeur.
\end{definition}

\begin{proposition}[Offres d'équilibre]\label{prop:qm:games-auctions-and-information:bids}
Avec $n$ enchérisseurs dont les valeurs sont indépendantes et de fonction de répartition $F$ sur $[0, 1]$~: dans l'\omterm{def:qm:games-auctions-and-information:auctions}{enchère au second prix}, enchérir sa propre valeur est une stratégie faiblement dominante~; dans l'\omterm{def:qm:games-auctions-and-information:auctions}{enchère au premier prix}, l'\omterm{def:qm:games-auctions-and-information:bayes}{équilibre de Bayes--Nash} symétrique
est $b(v) = v - \int_0^vF(x)^{n-1}dx/F(v)^{n-1}$, soit, pour des valeurs uniformes, $b(v) = \frac{n-1}nv$.
\end{proposition}

\begin{proof}
Second prix~: enchérir $b \ne v$ ne modifie le résultat que lorsque l'offre rivale la plus élevée $h$ se situe entre $b$ et $v$~; si $b > h > v$, l'enchérisseur gagne et paie $h > v$, une perte, et si $v > h > b$ il renonce au
profit $v - h$. Premier prix~: si les rivaux utilisent une fonction $b$ croissante, un enchérisseur de valeur $v$ qui offre $b(w)$ gagne avec probabilité $F(w)^{n-1}$ et réalise $(v - b(w))F(w)^{n-1}$. La condition du premier ordre en $w = v$ est
$\frac{d}{dv}(b(v)F(v)^{n-1}) = v\frac{d}{dv}F(v)^{n-1}$, et une intégration par parties à partir de $b(0) = 0$ donne la formule.
\end{proof}

Avec cinq enchérisseurs et des valeurs uniformes, chacun réduit son offre exactement d'un cinquième. L'équilibre est, enchérisseur par enchérisseur, une meilleure réponse~: un enchérisseur de valeur 0.8 face à quatre rivaux qui offrent $\frac45$ de leurs valeurs
maximise son profit espéré pour une offre de 0.64 (la \cref{fig:qm:games-auctions-and-information:br}). La sous-enchère est le prix de ne pas payer la deuxième offre la plus élevée~: elle compense exactement l'écart espéré entre
sa propre valeur et la suivante, immédiatement inférieure.

\begin{omfigure}
\begin{tikzpicture}
\begin{axis}[width=10cm,height=4.8cm,xlabel={offre},ylabel={profit espéré},
  tick label style={font=\small},label style={font=\small},xmin=0,xmax=0.8,ymin=0,ymax=0.08,grid=major,grid style={gray!25},
  scaled y ticks=false,yticklabel style={/pgf/number format/fixed,/pgf/number format/precision=2},table/col sep=comma]
\addplot[omDef,very thick] table[x=bid,y=payoff]{figdata/methods/29-games-auctions-and-information/bestresponse.csv};
\draw[omThm,dashed,thick] (axis cs:0.64,0) -- (axis cs:0.64,0.08) node[pos=0.9,left,font=\footnotesize]{$b(0.8) = 0.64$};
\end{axis}
\end{tikzpicture}

\omcaption{Profit espéré $(v - b)\,\P(\text{gagne})$ d'un enchérisseur de valeur 0.8 dans une \omterm{def:qm:games-auctions-and-information:auctions}{enchère au premier prix} face à quatre rivaux aux valeurs uniformes qui offrent $\frac45$ de leurs valeurs~: la meilleure réponse est
l'offre d'équilibre. Données~: le tutoriel du chapitre.}\label{fig:qm:games-auctions-and-information:br}
\end{omfigure}

\section{Formats d'enchères et recettes}

\begin{definition}[Enchères anglaise et hollandaise, prix de réserve]\label{def:qm:games-auctions-and-information:open}
Une \emph{enchère anglaise}\index{enchère anglaise} fait monter le prix jusqu'à ce qu'il ne reste qu'un enchérisseur, qui paie le prix auquel le dernier rival s'est retiré. Une \emph{enchère hollandaise}\index{enchère hollandaise} fait baisser le prix
jusqu'à ce qu'un enchérisseur l'accepte. Un \emph{prix de réserve}\index{prix de réserve} est le prix le plus bas auquel le vendeur accepte de vendre.
\end{definition}

L'\omterm{def:qm:games-auctions-and-information:open}{enchère anglaise} équivaut stratégiquement à une \omterm{def:qm:games-auctions-and-information:auctions}{enchère au second prix} à valeurs privées (rester en lice tant que le prix n'atteint pas sa valeur), et l'\omterm{def:qm:games-auctions-and-information:open}{enchère hollandaise} à une \omterm{def:qm:games-auctions-and-information:auctions}{enchère au premier prix} (choisir le prix auquel s'arrêter).

\begin{theorem}[Équivalence des recettes]\label{thm:qm:games-auctions-and-information:re}
Avec des valeurs privées indépendantes tirées d'une même loi continue et des enchérisseurs neutres au risque, toute enchère dans laquelle la valeur la plus élevée l'emporte à l'équilibre et un enchérisseur de la valeur la plus basse possible
espère un surplus nul dégage les mêmes recettes espérées (Vickrey, 1961~; Myerson, 1981~; Riley et Samuelson, 1981).
\end{theorem}

\begin{proof}[Idée]
À tout équilibre de ce type, le surplus espéré $U(v)$ d'un enchérisseur vérifie la condition d'enveloppe $U'(v) = \P(\text{gagne} \mid v) = F(v)^{n-1}$, donc $U(v) = U(0) + \int_0^vF^{n-1}$ est fixé par la règle d'allocation. Son
paiement espéré, $vF(v)^{n-1} - U(v)$, est alors le même dans chaque format, et il en va de même pour les recettes du vendeur, somme des paiements espérés.
\end{proof}

Avec cinq enchérisseurs et des valeurs uniformes, les recettes espérées de chaque format valent $\frac{n-1}{n+1} = \frac23$. Un million d'enchères simulées par format donnent 0.6665 (premier prix), 0.6665 (second prix), 0.6715 pour
une horloge anglaise qui monte par pas de cotation de 0.01 et 0.6614 pour une horloge hollandaise qui descend par pas de 0.01, les deux dernières s'écartant d'un demi-pas de cotation dans des sens opposés. Le gagnant au premier prix paie son offre sous-enchérie, le
gagnant au second prix paie la deuxième valeur~; les moyennes concordent, les distributions ne concordent pas.

\begin{proposition}[Prix de réserve optimal]\label{prop:qm:games-auctions-and-information:reserve}
Pour des lois de valeurs régulières, l'enchère maximisant les recettes fixe un \omterm{def:qm:games-auctions-and-information:open}{prix de réserve} $r^*$ solution de $r - (1 - F(r))/f(r) = v_0$, la propre valeur du vendeur, indépendant du nombre d'enchérisseurs (Myerson, 1981~;
admis). Pour des valeurs uniformes et $v_0 = 0$, $r^* = \frac12$, et avec $n$ enchérisseurs, les recettes espérées valent $\frac{n-1}{n+1} - \frac{2n}{n+1}r^{n+1} + r^n$.
\end{proposition}

Un \omterm{def:qm:games-auctions-and-information:open}{prix de réserve} accroît les recettes en renonçant parfois à vendre. Avec deux enchérisseurs, la réserve optimale de 0.5 fait passer les recettes espérées de $\frac13$ à $\frac5{12}$, soit une hausse de 25~\%, et laisse l'objet invendu un quart du temps~;
elle fait passer les recettes de 0.6667 à 0.6719 avec cinq enchérisseurs, soit 0.8~\%, et laisse l'objet invendu 3.1~\% du temps (la \cref{fig:qm:games-auctions-and-information:reserve})~: la concurrence fait le travail de la réserve.

\begin{omfigure}
\begin{tikzpicture}
\begin{axis}[width=10cm,height=5cm,xlabel={prix de réserve $r$},ylabel={recettes espérées},
  tick label style={font=\small},label style={font=\small},xmin=0,xmax=1,ymin=0,ymax=0.75,grid=major,grid style={gray!25},
  legend cell align=left,legend style={font=\footnotesize,at={(0.03,0.03)},anchor=south west,fill=white,draw=none},table/col sep=comma]
\addplot[omThm,very thick] table[x=r,y=n2]{figdata/methods/29-games-auctions-and-information/reserve.csv};
\addplot[omDef,very thick] table[x=r,y=n5]{figdata/methods/29-games-auctions-and-information/reserve.csv};
\draw[black,dashed] (axis cs:0.5,0) -- (axis cs:0.5,0.75);
\legend{deux enchérisseurs,cinq enchérisseurs}
\end{axis}
\end{tikzpicture}

\omcaption{Recettes espérées d'une enchère standard à valeurs privées uniformes en fonction du \omterm{def:qm:games-auctions-and-information:open}{prix de réserve}~; les deux courbes atteignent leur maximum en $r^* = \frac12$. Données~: le tutoriel du chapitre (forme close, vérifiée par
simulation).}\label{fig:qm:games-auctions-and-information:reserve}
\end{omfigure}

\begin{definition}[Enchères au prix offert et discriminatoires]\label{def:qm:games-auctions-and-information:pab}
Une \emph{enchère au prix offert}\index{enchère au prix offert}, ou \emph{enchère discriminatoire}\index{enchère discriminatoire}, adjuge plusieurs unités aux offres les plus élevées, chaque gagnant payant sa propre offre~; c'est la version à unités multiples
de l'\omterm{def:qm:games-auctions-and-information:auctions}{enchère au premier prix}, comme l'adjudication à prix uniforme du livre~2 est l'\omterm{def:qm:games-auctions-and-information:auctions}{enchère au second prix} à unités multiples.
\end{definition}

Le basculement du Trésor peut maintenant être modélisé. Prenons cinq dealers, chacun voulant l'une de deux unités, avec des valeurs privées uniformes. Dans l'adjudication à prix uniforme, chacun offre sa valeur, et les recettes valent le double de la troisième plus haute
valeur, $2 \times \frac36 = 1$ en espérance. Dans l'\omterm{def:qm:games-auctions-and-information:pab}{enchère au prix offert}, un dealer de valeur $v$ offre l'espérance de la deuxième valeur la plus élevée parmi les quatre autres, sachant qu'elle est inférieure à $v$, $b(v) = v(3 - 2.4v)/(4 -
3v)$~: 0.36 pour une valeur de 0.5, 0.54 pour 0.8, 0.60 pour 1. Par équivalence des recettes, les recettes espérées valent à nouveau 1~; 400\,000 enchères simulées donnent 0.9994 et 0.9999. Ce qui change, c'est tout le reste~: dans l'adjudication à prix uniforme, les
offres sont deux fois plus dispersées (l'écart entre les deux offres gagnantes vaut en moyenne 0.167 contre 0.087) et les recettes varient plus du double d'une enchère à l'autre (écart type 0.378 contre 0.160~;
la \cref{fig:qm:games-auctions-and-information:multi}), les deux traits que l'étude du Trésor a relevés. Avec une demande d'une seule unité, le modèle ne peut favoriser en moyenne aucun des deux formats~; les dealers réels demandent plusieurs unités et peuvent
pratiquer une \omterm{def:qm:games-auctions-and-information:auctions}{sous-enchère stratégique} (bid shading) dans une adjudication à prix uniforme (réduction de la demande), ce qui explique que la question soit empirique.

\begin{omfigure}
\begin{tikzpicture}
\begin{axis}[width=10cm,height=5cm,xlabel={recettes d'une enchère},ylabel={densité},
  tick label style={font=\small},label style={font=\small},xmin=0,xmax=2,ymin=0,ymax=3.8,grid=major,grid style={gray!25},
  legend cell align=left,legend style={font=\footnotesize,at={(0.97,0.97)},anchor=north east,fill=white,draw=none},table/col sep=comma]
\addplot[omThm,very thick] table[x=revenue,y=uniform]{figdata/methods/29-games-auctions-and-information/multiunit.csv};
\addplot[omDef,very thick] table[x=revenue,y=payasbid]{figdata/methods/29-games-auctions-and-information/multiunit.csv};
\legend{prix uniforme,au prix offert}
\end{axis}
\end{tikzpicture}

\omcaption{Distribution des recettes sur 400\,000 enchères simulées portant sur deux unités, avec cinq enchérisseurs à demande d'une seule unité et valeurs privées uniformes, selon les formats à prix uniforme et au prix offert à l'équilibre~: la même
moyenne, 1, et dispersions très différentes. Données~: le tutoriel du chapitre.}\label{fig:qm:games-auctions-and-information:multi}
\end{omfigure}

Les valeurs communes changent la logique. Si la valeur $V$ est commune et que chacun des cinq enchérisseurs observe $V$ plus un bruit uniforme sur $\pm0.1$, le signal le plus élevé surestime $V$ de $0.1 \times \frac46 = 0.067$ en moyenne, et un enchérisseur
qui offre son signal dans une \omterm{def:qm:games-auctions-and-information:auctions}{enchère au premier prix} perd autant lorsqu'il gagne, dans 97~\% de ses victoires~: la malédiction du vainqueur du livre~2. La sous-enchère fondée sur la surestimation espérée ramène le profit des gagnants à zéro. Avec
des valeurs affiliées, les formats ne rapportent plus les mêmes recettes~: l'\omterm{def:qm:games-auctions-and-information:open}{enchère anglaise}, qui révèle l'information des enchérisseurs éliminés, rapporte au moins autant que l'\omterm{def:qm:games-auctions-and-information:auctions}{enchère au second prix}, qui rapporte au moins autant que
l'\omterm{def:qm:games-auctions-and-information:auctions}{enchère au premier prix} (Milgrom et Weber, 1982).

\section{Entropie et pari optimal}

\begin{definition}[Entropie, information mutuelle]\label{def:qm:games-auctions-and-information:entropy}
L'\emph{entropie}\index{entropie} d'une variable aléatoire discrète $X$ est $H(X) = -\sum_xp(x)\log_2p(x)$ bits (Shannon, 1948). L'\emph{information mutuelle}\index{information mutuelle} de $X$ et $Y$ est $I(X; Y) =
\sum_{x,y}p(x, y)\log_2\frac{p(x, y)}{p(x)p(y)} = H(X) - H(X \mid Y)$, la \omterm{def:qm:estimation:kl}{divergence de Kullback--Leibler} de la loi conjointe au produit de ses lois marginales.
\end{definition}

\begin{theorem}[Kelly~: la valeur de l'information auxiliaire]\label{thm:qm:games-auctions-and-information:kelly}
Dans une course de chevaux aux probabilités de victoire $p(x)$ et aux cotes décimales $o(x)$, un parieur qui mise l'intégralité de sa richesse dans les proportions $b(x)$ a un taux de doublement $W(b) = \sum_xp(x)\log_2(b(x)o(x))$, maximisé pour $b = p$ (le
critère de Kelly du livre~2). Avec un signal $Y$ observé avant chaque course, le taux de doublement optimal augmente exactement de $I(X; Y)$ bits par course, quelles que soient les cotes (Kelly, 1956).
\end{theorem}

\begin{proof}
$W(b) - W(p) = \sum_xp(x)\log_2(b(x)/p(x)) = -D(p\,\|\,b) \le 0$, donc $b = p$ est optimal et $W^* = \sum_xp(x)\log_2(p(x)o(x))$. Avec le signal, les mises optimales sont $p(x \mid y)$ et
\[W^*_Y = \sum_{x,y}p(x, y)\log_2(p(x \mid y)o(x)).\] Les cotes s'annulent dans la différence~: $W^*_Y - W^* = \sum_{x,y}p(x, y)\log_2\frac{p(x \mid y)}{p(x)} = I(X; Y)$.
\end{proof}

Une course à quatre chevaux, de probabilités $(0.4, 0.3, 0.2, 0.1)$ et aux cotes décimales $(2.2, 3, 4.5, 9)$, où l'hippodrome prélève environ 11~\% (les probabilités implicites somment à 1.12), a une \omterm{def:qm:games-auctions-and-information:entropy}{entropie} de 1.846 bits. Le meilleur parieur
perd 0.165 bits par course~: sa richesse est divisée par deux toutes les six courses environ. Un pronostic qui désigne le gagnant avec une probabilité de 0.6 (et l'un des autres au hasard sinon) apporte 0.363 bits, et porte le taux de doublement à $+0.198$~;
200\,000 courses simulées donnent un gain de $0.363 \pm 0.002$. Un pronostic juste avec une probabilité $\frac14$ n'apporte rien, et un pronostic parfait apporte toute l'\omterm{def:qm:games-auctions-and-information:entropy}{entropie} (la \cref{fig:qm:games-auctions-and-information:kelly}).
L'information a un prix en bits, et les bits sont exactement la croissance.

\begin{omfigure}
\begin{tikzpicture}
\begin{axis}[width=10cm,height=5cm,xlabel={probabilité que le pronostic désigne le gagnant},ylabel={bits par course},
  tick label style={font=\small},label style={font=\small},xmin=0.2,xmax=1.02,ymin=0,ymax=2,grid=major,grid style={gray!25},
  legend cell align=left,legend style={font=\footnotesize,at={(0.03,0.97)},anchor=north west,fill=white,draw=none},table/col sep=comma]
\addplot[omDef,very thick] table[x=q,y=mi]{figdata/methods/29-games-auctions-and-information/kelly.csv};
\addplot[only marks,mark=*,mark size=1.6pt,omThm] table[x=q,y=sim]{figdata/methods/29-games-auctions-and-information/kelly.csv};
\legend{information mutuelle $I(X; Y)$,gain simulé du taux de doublement}
\end{axis}
\end{tikzpicture}

\omcaption{Gain du taux de doublement du parieur de Kelly grâce à un pronostic bruité, en fonction de la justesse du pronostic, dans une course de quatre chevaux~: l'\omterm{def:qm:games-auctions-and-information:entropy}{information mutuelle} et le gain mesuré sur 200\,000 courses simulées pour chaque justesse.
Données~: le tutoriel du chapitre, à graine fixée.}\label{fig:qm:games-auctions-and-information:kelly}
\end{omfigure}

\section{Tutoriel~: le basculement du Trésor}\label{tut:qm:games-auctions-and-information:switch}

\begin{tutorial}
\textbf{Objectif.} Vérifier l'équivalence des recettes dans quatre formats à unité unique et deux formats à unités multiples, fixer un \omterm{def:qm:games-auctions-and-information:open}{prix de réserve}, mettre en évidence la malédiction du vainqueur, et mesurer la valeur d'un pronostic en bits. \textbf{État final~:} les
figures et les nombres nommés du problème du week-end.
\begin{tutsteps}
\item \textbf{Offres d'équilibre.} \omterm{def:qm:games-auctions-and-information:auctions}{Enchère au premier prix} (forme close et forme générale), et \omterm{def:qm:games-auctions-and-information:pab}{enchère au prix offert} avec une demande d'une seule unité.
\omcode{code/firm/bidding/firm_bidding.py}{26}{58}{Fonctions d'offre d'équilibre (Python).}
\item \textbf{La valeur d'un pronostic.} La loi conjointe du gagnant et du pronostic, l'\omterm{def:qm:games-auctions-and-information:entropy}{information mutuelle} et les deux taux de doublement.
\omcode{code/methods/29-games-auctions-and-information/python/qm_games.py}{116}{139}{Kelly avec et sans information auxiliaire (Python).}
\item \textbf{Exécuter} \texttt{multiplicative\_weights(RPS)}, \texttt{best\_response\_curve()}, \texttt{equivalence()}, \texttt{reserve\_gain(2)}, \texttt{reserve\_gain(5)}, \texttt{treasury\_switch()},
\texttt{winners\_curse()} et \texttt{simulate\_races(0.6)} dans \texttt{qm\_games.py}~; puis \texttt{fig\_games.py}.
\end{tutsteps}
\textbf{Pour aller plus loin.} Donner à chaque dealer une demande de deux unités et trouver les offres à prix uniforme par itération de meilleure réponse~; rendre les enchérisseurs averses au risque et observer les recettes du premier prix dépasser celles du second.
\end{tutorial}

\section{Construction~: le simulateur d'enchères}\label{bld:qm:games-auctions-and-information:bidding}

\begin{build}
\bfield{Rôle} Un banc d'essai pour les offres de la firme dans les adjudications primaires, les transactions de bloc et les mises en concurrence par demande de cotation (RFQ)~: formats, indices de référence d'équilibre, \omterm{def:qm:games-auctions-and-information:open}{prix de réserve}.
\bfield{Interface} \texttt{fp\_bid}, \texttt{fp\_bid\_general}, \texttt{pab\_bid}~; \texttt{simulate(fmt, n, n\_auctions, seed, reserve, increment, values)} pour \texttt{first}, \texttt{second},
\texttt{english}, \texttt{dutch}~; \texttt{simulate\_multiunit(fmt, n, k, n\_auctions, seed)} pour \texttt{uniform}, \texttt{payasbid}~; \texttt{expected\_revenue\_uniform}~; \texttt{myerson\_reserve}~; \texttt{common\_value}.
\bfield{Règles} Chaque format simulé est vérifié par rapport à son indice de référence d'équilibre avant d'être utilisé pour évaluer une règle d'enchère~; les expériences à valeur commune rapportent toujours le profit des gagnants, et non celui des enchérisseurs.
\bfield{Tests d'acceptation} \texttt{code/firm/bidding/tests/}~: fonctions d'offre face aux formes closes (la formule générale du premier prix, l'\omterm{def:qm:games-auctions-and-information:pab}{enchère au prix offert} à une seule unité, égale au premier prix)~; équivalence des recettes dans quatre
formats~; \omterm{def:qm:games-auctions-and-information:open}{prix de réserve} pour des valeurs uniformes et non uniformes ainsi qu'une valeur positive pour le vendeur~; recettes et dispersions à unités multiples~; la malédiction du vainqueur.
\bfield{Pour aller plus loin} Demande à unités multiples et réduction de la demande~; enchérisseurs asymétriques par résolution numérique des équations différentielles d'équilibre~; signaux affiliés.
\end{build}

\begin{omsources}
\item P.~F.~Malvey et C.~M.~Archibald, ``Uniform-price auctions: update of the Treasury experience'', US Treasury, Office of Market Finance, octobre 1998~; TreasuryDirect, ``How auctions work''.
\item J.~F.~Nash, \emph{Proceedings of the National Academy of Sciences} 36, 1950~; J.~von Neumann, \emph{Mathematische Annalen} 100, 1928~; J.~C.~Harsanyi, \emph{Management Science} 14, 1967--68.
\item W.~Vickrey, ``Counterspeculation, auctions, and competitive sealed tenders'', \emph{Journal of Finance} 16, 1961~; R.~B.~Myerson, ``Optimal auction design'', \emph{Mathematics of Operations Research} 6,
1981~; J.~G.~Riley et W.~F.~Samuelson, ``Optimal auctions'', \emph{American Economic Review} 71, 1981~; P.~R.~Milgrom et R.~J.~Weber, \emph{Econometrica} 50, 1982.
\item C.~E.~Shannon, \emph{Bell System Technical Journal} 27, 1948~; J.~L.~Kelly, ``A new interpretation of information rate'', \emph{Bell System Technical Journal} 35, 1956~; T.~M.~Cover et J.~A.~Thomas,
\emph{Elements of Information Theory}, 2\up{e} éd., Wiley, 2006.
\end{omsources}

\section{Exercices}

\begin{exercise}[$\star$]\label{exo:qm:games-auctions-and-information:1}
Dix enchérisseurs aux valeurs privées uniformes enchérissent dans une \omterm{def:qm:games-auctions-and-information:auctions}{enchère au premier prix}. Que mise un enchérisseur de valeur 0.7, et quel est le revenu espéré~?
\end{exercise}

\begin{exercise}[$\star$]\label{exo:qm:games-auctions-and-information:2}
Déterminer l'équilibre mixte du jeu matching pennies de gains $\pm1$, ainsi que sa valeur.
\end{exercise}

\begin{exercise}[$\star$]\label{exo:qm:games-auctions-and-information:3}
Calculer l'\omterm{def:qm:games-auctions-and-information:entropy}{entropie} d'un dé équilibré, ainsi que l'\omterm{def:qm:games-auctions-and-information:entropy}{information mutuelle} entre le dé et sa parité.
\end{exercise}

\begin{exercise}[$\star\star$]\label{exo:qm:games-auctions-and-information:4}
Déterminer l'offre d'équilibre au premier prix pour des valeurs de fonction de répartition $F(v) = v^2$ sur $[0, 1]$, puis la vérifier à l'aide de la formule générale.
\end{exercise}

\begin{exercise}[$\star\star$]\label{exo:qm:games-auctions-and-information:5}
Déterminer le \omterm{def:qm:games-auctions-and-information:open}{prix de réserve} optimal pour des valeurs uniformes sur $[0, 1]$ lorsque le vendeur évalue l'objet à 0.2.
\end{exercise}

\begin{exercise}[$\star\star$]\label{exo:qm:games-auctions-and-information:6}
Montrer que, dans une \omterm{def:qm:games-auctions-and-information:pab}{enchère au prix offert} de $k$ unités adressée à des enchérisseurs à demande unitaire, le cas $k = 1$ redonne l'équilibre de l'\omterm{def:qm:games-auctions-and-information:auctions}{enchère au premier prix}.
\end{exercise}

\begin{exercise}[$\star\star\star$]\label{exo:qm:games-auctions-and-information:7}
\emph{Programmation.} Simuler les quatre formats à unité unique avec cinq enchérisseurs, un million d'enchères pour chacun, avec et sans \omterm{def:qm:games-auctions-and-information:open}{prix de réserve} optimal, et comparer les résultats à la forme fermée.
\end{exercise}

\begin{exercise}[$\star\star\star$]\label{exo:qm:games-auctions-and-information:8}
\emph{Repérer le défaut.} «~Nous avons remporté les 30 dernières transactions de bloc sur lesquelles nous avons soumis une offre, chacune au fair value de notre modèle~; le modèle est excellent.»
\end{exercise}

\section{Problème~: le changement de cap du Trésor}

\begin{problem}[{Problème du week-end --- payer son offre ou payer le prix d'équilibre}]\label{pb:qm:games-auctions-and-information:1}
Cinq dealers enchérissent sur des titres~; le Trésor se demande si un prix d'équilibre unique rapporterait davantage qu'une \omterm{def:qm:games-auctions-and-information:pab}{enchère au prix offert}.

\medskip\noindent\textbf{Partie I --- Équilibre.}
\begin{enumerate}
\item Quelle offre chaque dealer formule-t-il dans une \omterm{def:qm:games-auctions-and-information:auctions}{enchère au premier prix} à unité unique et à valeurs uniformes, et pourquoi~?
\item Pourquoi l'offre sincère est-elle dominante dans l'\omterm{def:qm:games-auctions-and-information:auctions}{enchère au second prix}~?
\item Que montre la courbe de meilleure réponse d'un dealer de valeur 0.8~?
\item Que prédit l'équivalence des revenus pour les quatre formats à unité unique, et que donne la simulation~?
\item Pourquoi les horloges des \omterm{def:qm:games-auctions-and-information:open}{enchères anglaise} et hollandaise, avec des pas de cotation de 0.01, manquent-elles d'un demi-pas dans des sens opposés~?
\end{enumerate}

\medskip\noindent\textbf{Partie II --- \omterm{def:qm:games-auctions-and-information:open}{Prix de réserve} et enchères multi-unités.}
\begin{enumerate}[resume]
\item Quel est le \omterm{def:qm:games-auctions-and-information:open}{prix de réserve} optimal et quel gain procure-t-il avec deux puis avec cinq enchérisseurs~?
\item Que proposent les dealers dans l'\omterm{def:qm:games-auctions-and-information:pab}{enchère au prix offert} à deux unités~?
\item Quel revenu rapportent les deux formats multi-unités~?
\item Qu'est-ce qui les distingue, et comment cela correspond-il à l'étude du Trésor~?
\item Pourquoi le modèle ne pouvait-il pas trancher la question posée par le Trésor~?
\end{enumerate}

\medskip\noindent\textbf{Partie III --- Information.}
\begin{enumerate}[resume]
\item Quelle est l'ampleur de la malédiction du vainqueur dans l'expérience à valeur commune, et comment la supprime-t-on~?
\item Quelle est l'\omterm{def:qm:games-auctions-and-information:entropy}{entropie} de la course, et quel est le taux de doublement sans pronostic~?
\item Que rapporte un pronostic juste dans 60~\% des cas, et que donne la simulation~?
\item Pourquoi la cote n'influe-t-elle pas sur le gain~?
\item Que vaut un pronostic juste dans 25~\% des cas, et pourquoi~?
\end{enumerate}

\medskip\noindent\textbf{Partie IV --- Jugement.}
\begin{enumerate}[resume]
\item Qu'est-ce qui ferait rapporter davantage, en pratique, l'adjudication à prix uniforme~?
\item Quand la firme doit-elle pratiquer une \omterm{def:qm:games-auctions-and-information:auctions}{sous-enchère stratégique} (bid shading) dans une compétition pour des transactions de bloc~?
\item Combien est-on prêt à payer pour un pronostic, en taux de doublement~?
\item Énoncer le \emph{résultat nommé}~: le revenu espéré des deux formats dans le modèle symétrique à valeurs privées, la sous-enchère d'équilibre de cinq dealers, et le gain lié au \omterm{def:qm:games-auctions-and-information:open}{prix de réserve} optimal.
\item En une phrase~: que dit l'équivalence des revenus sur la conception des enchères~?
\end{enumerate}
\end{problem}

\section{Questions d'entretien}

\begin{interviewq}[$\star$ \iqroles{trader, chercheur}]\label{iq:qm:games-auctions-and-information:1}
Deux enchérisseurs aux valeurs uniformes sur $[0, 1]$ enchérissent dans une \omterm{def:qm:games-auctions-and-information:auctions}{enchère au premier prix}. Quelle offre devriez-vous formuler avec la valeur $v$~?
\end{interviewq}

\begin{interviewq}[$\star\star$ \iqroles{trader}]\label{iq:qm:games-auctions-and-information:2}
Qu'est-ce que la malédiction du vainqueur et comment devez-vous enchérir pour l'éviter~?
\end{interviewq}

\begin{interviewq}[$\star\star$ \iqroles{chercheur}]\label{iq:qm:games-auctions-and-information:3}
Énoncez le théorème d'équivalence des revenus et donnez deux façons dont il échoue en pratique.
\end{interviewq}

\begin{interviewq}[$\star\star$ \iqroles{chercheur, trader}]\label{iq:qm:games-auctions-and-information:4}
Quelle est la valeur d'un signal pour un parieur suivant le critère de Kelly~?
\end{interviewq}

\begin{interviewq}[$\star\star$ \iqroles{chercheur}]\label{iq:qm:games-auctions-and-information:5}
Déterminez la valeur et les stratégies optimales du \omterm{def:qm:games-auctions-and-information:game}{jeu à somme nulle} de matrice des gains $\begin{pmatrix}3 & -1\\ -2 & 1\end{pmatrix}$ pour le joueur de ligne.
\end{interviewq}

\begin{interviewq}[$\star\star\star$ \iqroles{trader, chercheur}]\label{iq:qm:games-auctions-and-information:6}
Vous soumettez une offre sur un portefeuille dans une demande de cotation (RFQ) face à quatre dealers. Comment le nombre de concurrents modifie-t-il votre offre~?
\end{interviewq}
